\documentclass[10pt,a4paper,landscape,utf8]{ctexart}

% ====== 极致压缩页边距 ======
\usepackage[left=0.5cm,right=0.5cm,top=0.6cm,bottom=0.6cm]{geometry}
\usepackage{multicol}
\usepackage{amsmath, amssymb, bm}
\usepackage{esint} 

% ====== 自定义紧凑标题 ======
\usepackage{titlesec}
\titleformat{\section}{\normalfont\bfseries\small}{\thesection}{1em}{}
\titlespacing*{\section}{0pt}{1.5ex plus .2ex}{0.5ex plus .1ex}

% ====== 宏定义 ======
\newcommand{\p}{\partial}
\newcommand{\bs}[1]{\boldsymbol{#1}}
\newcommand{\mr}[1]{\mathrm{#1}}

\begin{document}
\scriptsize % 全局小字号
\setlength{\parindent}{0pt}
\setlength{\columnsep}{0.6cm}
\setlength{\columnseprule}{0.4pt}

\begin{multicols*}{3}

% ==========================================
$$ \epsilon_{ijk}\epsilon_{lmn} = \begin{vmatrix} \delta_{il} & \delta_{im} & \delta_{in} \\ \delta_{jl} & \delta_{jm} & \delta_{jn} \\ \delta_{kl} & \delta_{km} & \delta_{kn} \end{vmatrix} $$
$$ \sum_{k} \epsilon_{ijk}\epsilon_{kmn} = \delta_{im}\delta_{jn} - \delta_{in}\delta_{jm} $$
$$ \sum_{j,k} \epsilon_{ijk}\epsilon_{mjk} = 2\delta_{im}, \quad \sum_{i,j,k} \epsilon_{ijk}\epsilon_{ijk} = 3! = 6 $$
$$ \bs{A}\times(\bs{B}\times\bs{C}) = \bs{B}(\bs{A}\cdot\bs{C}) - \bs{C}(\bs{A}\cdot\bs{B}) $$
$$ \nabla(\bs{A}\cdot\bs{B}) = \bs{A}\times(\nabla\times\bs{B}) + \bs{B}\times(\nabla\times\bs{A}) + (\bs{A}\cdot\nabla)\bs{B} + (\bs{B}\cdot\nabla)\bs{A} $$
$$ \nabla \cdot (\bs{A} \times \bs{B}) = \bs{B} \cdot (\nabla \times \bs{A}) - \bs{A} \cdot (\nabla \times \bs{B}) $$
$$ \nabla \times (\nabla \times \bs{A}) = \nabla(\nabla \cdot \bs{A}) - \nabla^2\bs{A} $$
\textbf{相对位置矢量 $\bs{R} = \bs{r} - \bs{r}'$ :} $R = |\bs{r} - \bs{r}'|$
$$ \nabla \cdot \bs{R} = 3, \quad \nabla \times \bs{R} = 0 $$
$$ \nabla R = \frac{\bs{R}}{R} = \hat{R}, \quad \nabla \left(\frac{1}{R}\right) = -\frac{\hat{R}}{R^2} $$
$$ \nabla \cdot \left(\frac{\hat{R}}{R^2}\right) = 4\pi\delta^3(\bs{R}) \implies \nabla^2 \left(\frac{1}{R}\right) = -4\pi\delta^3(\bs{r} - \bs{r}') $$

% ==========================================

静电力 $\bs{F} = q\bs{E}$, 电场 $\bs{E} = -\nabla V$。

\textbf{基本方程:} $\nabla \times \bs{E} = 0 \implies \bs{E} = -\nabla V$.

高斯定律积分形式：$$ \oint \bs{E} \cdot d\bs{A} = \frac{Q_{\text{enc}}}{\epsilon_0} $$

泊松方程: $\nabla^2 V = -\frac{\rho}{\epsilon_0}$. 拉普拉斯方程: $\nabla^2 V = 0$.

\textbf{直接积分法 (库仑定律):}
$$ V(\bs{r}) = \frac{1}{4\pi\epsilon_0} \int \frac{\rho(\bs{r}')}{|\bs{r}-\bs{r}'|} d\tau' $$
* \textbf{方形线圈} (边长 $a$，求 $z$ 轴上的点): 
$$V = \frac{\lambda}{\pi\epsilon_0}\ln\left(\frac{\sqrt{z^2+a^2/2}+a/2}{\sqrt{z^2+a^2/2}-a/2}\right)$$
* \textbf{圆形线圈} (半径 $r$): $V = \frac{Q}{4\pi\epsilon_0\sqrt{r^2+z^2}}$

* \textbf{扁平面圆盘} (半径 $R$): $V = \frac{\sigma}{2\epsilon_0}(\sqrt{R^2+z^2} - |z|)$

\textbf{静电场能量:}
$$ W = \frac{1}{2} \int \rho V d\tau = \frac{\epsilon_0}{2} \int E^2 d\tau $$

% ==========================================

\textbf{位于接地平面上方 $z=d$ 处的点电荷 $q$:}

* 镜像电荷: 位于 $z=-d$ 处的 $-q$。

* 受力: $\bs{F} = -\frac{q^2}{4\pi\epsilon_0 (2d)^2}\hat{z}$

* 诱导表面电荷密度: $\sigma(x,y) = -\frac{qd}{2\pi(x^2+y^2+d^2)^{3/2}}$

* 系统总能量: $W = -\frac{q^2}{16\pi\epsilon_0 d}$ (真实电荷与镜像电荷系统能量的一半！)。

\textbf{相交接地平面 (夹角 $\alpha$):}需要的镜像电荷数量: $N = \frac{360^\circ}{\alpha} - 1$。

* 仅当 $180^\circ/\alpha \in \mathbb{Z}^+$ 时有效 (例如 $90^\circ, 60^\circ, 45^\circ$)。

静电学界面条件：$$E_{out}-E_{in} = \frac{\sigma}{\epsilon_0}$$

% ==========================================

当 $r \gg r'$ 时，通过泰勒级数展开 $\frac{1}{|\bs{r}-\bs{r}'|}$ :
$$ V(\bs{r}) = \frac{1}{4\pi\epsilon_0} \sum_{n=0}^{\infty} \frac{1}{r^{n+1}} \int (r')^n P_n(\cos\theta') \rho(\bs{r}') d\tau' $$

球坐标系下具有方位角对称性（与 $\phi$ 无关）的拉普拉斯通解为：$$ V(r, \theta) = \sum_{l=0}^{\infty} \left( A_l r^l + \frac{B_l}{r^{l+1}} \right) P_l(\cos\theta) $$

$P_1(x)=x$, $P_2(x)=\frac{1}{2}(3x^2-1)$, $P_3(x)=\frac{1}{2}(5x^3-3x)$。

\textbf{单极子 (Monopole):} $V_{mon} = \frac{Q}{4\pi\epsilon_0 r}$ (其中 $Q = \int \rho d\tau'$)

\textbf{偶极矩 (Dipole Moment):} $\bs{p} = \int \bs{r}' \rho(\bs{r}') d\tau'$ (离散形式: $\sum q_i \bs{r}_i$)
$$ V_{dip}(\bs{r}) = \frac{1}{4\pi\epsilon_0} \frac{\bs{p} \cdot \hat{r}}{r^2} $$
$$ \bs{E}_{dip}(\bs{r}) = \frac{1}{4\pi\epsilon_0 r^3} [3(\bs{p}\cdot\hat{r})\hat{r} - \bs{p}] $$
* 注：若总电荷 $Q=0$，则 $\bs{p}$ 的值与坐标原点的选择无关。

% ==========================================

极化强度 $\bs{P}$ = 单位体积的偶极矩。

\textbf{束缚电荷 (Bound Charges):}
$$ \rho_b = -\nabla \cdot \bs{P} \quad \text{(体束缚电荷)} $$
$$ \sigma_b = \bs{P} \cdot \hat{n} \quad \text{(面束缚电荷)} $$
* $\hat{n}$ 是电介质表面的 \textbf{向外} 单位法向量。

* 总束缚电荷始终为零: $\int \rho_b d\tau + \oint \sigma_b da = 0$。

\textbf{电位移矢量 $\bs{D}$:}
$$ \bs{D} = \epsilon_0\bs{E} + \bs{P} $$
$\bs{D}$ 的高斯定理 (仅包含自由电荷):
$$ \nabla \cdot \bs{D} = \rho_f \implies \oint \bs{D} \cdot d\bs{A} = Q_{f, \text{enc}} $$
* 线性、各向同性、均匀 (LIH) 介质:
$$ \bs{P} = \epsilon_0 \chi_e \bs{E}, \quad \bs{D} = \epsilon_0(1+\chi_e)\bs{E} = \epsilon \bs{E} = \epsilon_r\epsilon_0\bs{E} $$

\textbf{边界条件 (一般形式):}
$$ \bs{E}_{1\parallel} - \bs{E}_{2\parallel} = 0 \implies \text{切向 } \bs{E} \text{ 连续。} $$
$$ \bs{D}_{1\perp} - \bs{D}_{2\perp} = \sigma_f \implies \text{法向 } \bs{D} \text{ 突变 } \sigma_f \text{。} $$
$$ V_1 = V_2 \quad \text{(电势始终连续)} $$
* 计算电介质中的电容套路: 由 $Q_f$ 求 $D$ $\to$ $E = D/\epsilon$ $\to$ $V = \int E dl$ $\to$ $C = Q_f / V$。

% ==========================================

\textbf{毕奥-萨伐尔定律 (Biot-Savart Law):}
$$ \bs{B}(\bs{r}) = \frac{\mu_0 I}{4\pi} \int \frac{d\bs{l}' \times \hat{R}}{R^2} $$
\textbf{安培定律 (Ampere's Law):} (适用于高对称性，如圆柱、平面)
$$ \oint \bs{B} \cdot d\bs{l} = \mu_0 I_{\text{enc}} \implies \nabla \times \bs{B} = \mu_0\bs{J} $$
* 长同轴电缆 ($a < s < b$ 间隙处): $B = \frac{\mu_0 I}{2\pi s}\hat{\phi}$

* 螺线管 (内部): $B = \mu_0 n I$ ($n$ = 单位长度匝数)

\textbf{磁矢势 $\bs{A}$ (Magnetic Vector Potential):}
$$ \bs{B} = \nabla \times \bs{A}, \quad \nabla \cdot \bs{B} = 0 $$

远场标准磁偶极子: $$\bs{A} = \frac{\mu_0}{4\pi} \frac{\bs{m} \times \hat{r}}{r^2} = \frac{\mu_0 (\pi \lambda \omega R^3)}{4\pi} \frac{\hat{z} \times \hat{r}}{r^2} = \frac{\mu_0 \lambda \omega R^3}{4 r^2} \sin\theta \hat{\phi}$$

$$\bs{B} = \frac{\mu_0 m}{4\pi r^3} (2\cos\theta \hat{r} + \sin\theta \hat{\theta}) = \frac{\mu_0 \lambda \omega R^3}{4 r^3} (2\cos\theta \hat{r} + \sin\theta \hat{\theta})$$

库仑规范 ($\nabla \cdot \bs{A} = 0$) $\implies \nabla^2\bs{A} = -\mu_0\bs{J}$
$$ \bs{A}(\bs{r}) = \frac{\mu_0}{4\pi} \int \frac{\bs{J}(\bs{r}')}{|\bs{r}-\bs{r}'|} d\tau' $$

\textbf{磁性介质与束缚电流:}
磁化强度 $\bs{M}$ (单位体积的磁偶极矩)。
$$ \bs{J}_b = \nabla \times \bs{M} \quad \text{(体电流)}, \quad \bs{K}_b = \bs{M} \times \hat{n} \quad \text{(面电流)} $$
$$ \bs{H} = \frac{\bs{B}}{\mu_0} - \bs{M} \implies \nabla \times \bs{H} = \bs{J}_f $$
线性介质: $\bs{M} = \chi_m \bs{H}$, $\bs{B} = \mu\bs{H} = \mu_r\mu_0\bs{H}$。

% ==========================================

\textbf{法拉第电磁感应定律:}
$$ \mathcal{E} = \oint \bs{E} \cdot d\bs{l} = -\frac{d\Phi_B}{dt} = -\frac{d}{dt} \int \bs{B} \cdot d\bs{A} $$
微分形式: $\nabla \times \bs{E} = -\frac{\p\bs{B}}{\p t}$.

* 楞次定律: 感应电流的方向总是阻碍磁通量 $\Phi_B$ 的变化。

* \textit{在半径为 $s$ 的圆形区域内有均匀变化的 $\bs{B}$:}
$$E(r<s) = -\frac{r}{2}\frac{dB}{dt}; \quad E(r>s) = -\frac{s^2}{2r}\frac{dB}{dt}$$

\textbf{电感 ($L$) 与 磁能 ($U$):}
$$ L = \frac{\Phi_B}{I}, \quad \mathcal{E} = -L\frac{dI}{dt} $$
磁能密度: $u_B = \frac{1}{2}\bs{B}\cdot\bs{H} = \frac{B^2}{2\mu_0}$
总磁能:
$$ U = \int u_B d\tau = \frac{1}{2}LI^2 \implies L = \frac{2U}{I^2} $$
* \textit{同轴电缆单位长度自感:} $L/l = \frac{\mu_0}{2\pi}\ln(b/a)$。

% ==========================================

\textbf{麦克斯韦方程组 (介质中):}

1. $\nabla \cdot \bs{D} = \rho_f$ (高斯定律)

2. $\nabla \cdot \bs{B} = 0$ (无磁单极子定律)

3. $\nabla \times \bs{E} = -\frac{\p\bs{B}}{\p t}$ (法拉第定律)

4. $\nabla \times \bs{H} = \bs{J}_f + \frac{\p\bs{D}}{\p t}$ (安培-麦克斯韦定律)

\textbf{坡印廷定理 (能量守恒):}
$$ \frac{\p u}{\p t} + \nabla \cdot \bs{S} = -\bs{J}_f \cdot \bs{E} $$
坡印廷矢量 (能流密度): $\bs{S} = \bs{E} \times \bs{H} = \frac{1}{\mu}(\bs{E} \times \bs{B})$。

% ==========================================

\textbf{波动方程 (无源、LIH 介质):}
$$ \nabla^2\bs{E} = \mu\epsilon\frac{\p^2\bs{E}}{\p t^2}, \quad v = \frac{1}{\sqrt{\mu\epsilon}} $$
\textbf{平面波:} $\bs{E} = \bs{E}_0 e^{i(\bs{k}\cdot\bs{r} - \omega t)}$

* 色散关系: $k = \omega\sqrt{\mu\epsilon} \implies \omega = vk$。

* 正交性: $\bs{k} \cdot \bs{E} = 0, \quad \bs{B} = \frac{1}{\omega}(\bs{k} \times \bs{E}) = \frac{1}{v}(\hat{k} \times \bs{E})$。

* 振幅关系: $B_0 = E_0 / v$。

* 能量: $u_E = u_B$, $u = \epsilon E^2$. 时间平均 $\langle\bs{S}\rangle = \frac{1}{2}\epsilon v E_0^2 \hat{k}$。

\textbf{复介电常数 (吸收介质):}
若 $\epsilon = \epsilon_1 + i\epsilon_2$, 波矢变为复数 $k = k_1 + ik_2$。
$$ k_{1} = \omega\sqrt{\frac{\mu_0(|\epsilon|+\epsilon_1)}{2}}, \quad k_{2} = \omega\sqrt{\frac{\mu_0(|\epsilon|-\epsilon_1)}{2}} $$
* 衰减: 波的振幅随 $e^{-k_2 z}$ 呈指数衰减。

* 趋肤深度: $\delta = 1/k_2$。相速度 $v_p = \omega/k_1$。

\textbf{波的叠加与速度:}
两列波 $\omega_1 \approx \omega_2$, $k_1 \approx k_2$ 叠加会产生拍频 (beats)。

* 载波相速度 (Phase velocity): $v_p = \bar{\omega}/\bar{k} = \omega/k$。

* 包络群速度 (Group velocity): $v_g = d\omega/dk$。

% ==========================================

界面位于 $x=0$。入射波 (I)、反射波 (R)、透射波 (T):
$$E_I = \tilde{E}_{0I} e^{i(k_{1}x\cos \theta_I + k_{1}y\sin \theta_I - \omega t)}\hat{z}$$
$$E_R = \tilde{E}_{0R} e^{i(-k_{1}x\cos \theta_R + k_{1}y\sin \theta_R - \omega t)}\hat{z}$$
$$E_T = \tilde{E}_{0T} e^{i(k_{2}x\cos \theta_T + k_{2}y\sin \theta_T - \omega t)}\hat{z}$$
斯涅尔定律 (Snell's Law): $\theta_I = \theta_R, \quad n_1\sin\theta_I = n_2\sin\theta_T$。

\textbf{边界条件:}
$\bs{E}_{\parallel}$，$\bs{B}_{\perp}$ 连续；$\bs{B}_{\parallel}/\mu$，$\bs{D}_{\perp}$ 连续。

\textbf{TE 模式 / s-极化 ($\bs{E} \perp$ 入射面):}
令 $\alpha = \frac{\cos\theta_T}{\cos\theta_I}$, $\beta = \frac{\mu_1 v_1}{\mu_2 v_2} \approx \frac{n_2}{n_1}$ (若 $\mu_1=\mu_2$)。
$$ r = \frac{\tilde{E}_{0R}}{\tilde{E}_{0I}} = \frac{1 - \alpha\beta}{1 + \alpha\beta}, \quad t = \frac{\tilde{E}_{0T}}{\tilde{E}_{0I}} = \frac{2}{1 + \alpha\beta} $$
* 反射率 $R = |r|^2$. 透射率 $T = \alpha\beta |t|^2$. $R+T=1$.

* \textit{布儒斯特角:} 发生于反射波振幅 $r=0$ 时。对于 TE 模式，要求 $1-\alpha\beta=0 \implies n_1=n_2$。因此，\textbf{TE 模式不存在布儒斯特角}。

\textbf{TM 模式 / p-极化 ($\bs{B} \perp$ 入射面):}
$$ r = \frac{\alpha - \beta}{\alpha + \beta}, \quad t = \frac{2}{\alpha + \beta} $$
* \textit{布儒斯特角:} 发生于 $\alpha=\beta$ 时。$\theta_B \approx \arctan(n_2/n_1)$。

% ==========================================

由 $\phi$ 和 $\bs{A}$ 求解电磁场:
$$ \bs{E} = -\nabla\phi - \frac{\p\bs{A}}{\p t}, \quad \bs{B} = \nabla \times \bs{A} $$
* 洛伦兹规范: $\nabla \cdot \bs{A} = -\mu_0\epsilon_0\frac{\p\phi}{\p t}$。这使得 $\phi$ 和 $\bs{A}$ 的波动方程完全解耦。

* \textit{示例:} 若 $\phi=0$, $\bs{A} = A_0\sin(kx-\omega t)\hat{y}$。

则 $\bs{E} = A_0\omega\cos(\dots)\hat{y}$, $\bs{B} = A_0 k\cos(\dots)\hat{z}$。

必须要求 $\omega = ck$ 才能满足真空麦克斯韦方程组。

% ==========================================

\textbf{球坐标系基矢转换:}
$x=r\sin\theta\cos\phi$, $y=r\sin\theta\sin\phi$, $z=r\cos\theta$.

$\hat{x}= \sin\theta\cos\phi\hat{r} + \cos\theta\cos\phi\hat{\theta} - \sin\phi\hat{\phi}$

$\hat{y}= \sin\theta\sin\phi\hat{r} + \cos\theta\sin\phi\hat{\theta} + \cos\phi\hat{\phi}$

$ \hat{z} = \cos\theta \hat{r} - \sin\theta \hat{\theta} $

$\hat{\phi} = -\sin\phi\hat{x} + \cos\phi\hat{y}$

$\hat{\theta} = \cos\theta\cos\phi\hat{x} + \cos\theta\sin\phi\hat{y} - \sin\theta\hat{z}$

$\hat{r} = \sin\theta\cos\phi\hat{x} + \sin\theta\sin\phi\hat{y} + \cos\theta\hat{z}$

\textbf{球坐标系正交关系:}

$\hat{r} \times \hat{\theta} = \hat{\phi}, \quad \hat{\theta} \times \hat{\phi} = \hat{r}, \quad \hat{\phi} \times \hat{r} = \hat{\theta}$

\textbf{球坐标系下的梯度与散度:}
$$ \nabla f = \frac{\p f}{\p r}\hat{r} + \frac{1}{r}\frac{\p f}{\p \theta}\hat{\theta} + \frac{1}{r\sin\theta}\frac{\p f}{\p \phi}\hat{\phi} $$
$$ \nabla \cdot \bs{A} = \frac{1}{r^2}\frac{\p}{\p r}(r^2 A_r) + \frac{1}{r\sin\theta}\frac{\p}{\p \theta}(\sin\theta A_\theta) + \frac{1}{r\sin\theta}\frac{\p A_\phi}{\p \phi} $$
% ==========================================

\textbf{1. 电介质与电容:}
$$ Q_f \xrightarrow{\oint \bs{D}\cdot d\bs{A} = Q_f} \bs{D} \xrightarrow{\bs{E}=\frac{\bs{D}}{\epsilon}} \bs{E} \xrightarrow{\int \bs{E}\cdot d\bs{l}} V \implies C = \frac{Q_f}{V} $$
* \textbf{求束缚电荷:} $\bs{P} = \bs{D} - \epsilon_0\bs{E} \implies \rho_b = -\nabla\cdot\bs{P}, \sigma_b = \bs{P}\cdot\hat{n}_{out}$

\textbf{2. 磁介质与电感:}
$$ I_f \xrightarrow{\oint \bs{H}\cdot d\bs{l} = I_f} \bs{H} \xrightarrow{\bs{B}=\mu\bs{H}} \bs{B} \xrightarrow{\int \bs{B}\cdot d\bs{A}} \Phi_B \implies L = \frac{\Phi_B}{I_f} $$
* \textbf{能量法求电感:} $U = \int \frac{B^2}{2\mu} d\tau \implies L = \frac{2U}{I_f^2}$

* \textbf{求束缚电流:} $\bs{M} = \frac{\bs{B}}{\mu_0} - \bs{H} \implies \bs{J}_b = \nabla\times\bs{M}, \bs{K}_b = \bs{M}\times\hat{n}_{out}$

\textbf{3. 变磁场与感应电场:}
$$ \bs{B}(t) \xrightarrow{\int \bs{B}\cdot d\bs{A}} \Phi_B(t) \xrightarrow{-\frac{d}{dt}} \mathcal{E} \xrightarrow{\oint \bs{E}\cdot d\bs{l} = \mathcal{E}} \bs{E}_{ind} $$
* \textbf{能量流验证:} $\bs{S} = \frac{1}{\mu_0}(\bs{E}\times\bs{B}) \xrightarrow{\oint \bs{S}\cdot d\bs{A}} P_{in} = \frac{dU_B}{dt}$

\textbf{4. 电磁波界面透射/反射:}

\textcircled{1} :画出界面，利用斯涅尔定律定 $\bs{k}_I, \bs{k}_R, \bs{k}_T$ 方向。\\
\textcircled{2} : $\bs{E} \perp$ 入射面(TE) 或 $\bs{B} \perp$ 入射面(TM)。\\
\textcircled{3} : 用 $\bs{B} = \frac{1}{v}(\hat{k}\times\bs{E})$ 或 $\bs{E} = -v(\hat{k}\times\bs{B})$ 补全所有场。\\
\textcircled{4} : 界面处 $\bs{E}_{\parallel}$ 连续，$\bs{B}_{\parallel}/\mu$ 连续。\\
\textcircled{5} : 联立方程组，解出 $r$ 和 $t$。

% ==========================================
* \textbf{Permittivity} - 介电常数 ($\epsilon$)

* \textbf{Permeability} - 磁导率 ($\mu$)

* \textbf{Method of Images} - 镜像法

* \textbf{Multipole Expansion} - 多极展开

* \textbf{Electric Displacement} - 电位移矢量 ($\bs{D}$)

* \textbf{Magnetic Vector Potential} - 磁矢势 ($\bs{A}$)

* \textbf{Poynting Vector} - 坡印廷矢量 / 能流密度 ($\bs{S}$)

* \textbf{Dispersion Relation} - 色散关系

* \textbf{Phase Velocity / Group Velocity} - 相速度 / 群速度

* \textbf{Attenuation} - 衰减

* \textbf{Brewster's Angle} - 布儒斯特角

* \textbf{Polarization} - 极化 (介质中) / 偏振 (电磁波)

* \textbf{Indutance} - 电感 ($L$)


设带电量为 $q$ 的点电荷位于坐标原点，一个半径为 $R$、长度为 $L$、均匀体电荷密度为 $\rho$ 的实心圆柱体被放置在 $z$ 轴上。圆柱靠近原点的一端位于 $z=d$ 处，较远的一端位于 $z=d+L$ 处。

$$ E_z = \int_d^{d+L} \frac{\rho}{2\epsilon_0} \left( 1 - \frac{z}{\sqrt{z^2 + R^2}} \right) dz $$
$$= \frac{\rho}{2\epsilon_0} \left( L + \sqrt{d^2 + R^2} - \sqrt{(d+L)^2 + R^2} \right) $$

$$ F = \frac{q \rho}{2\epsilon_0} \left( L + \sqrt{d^2 + R^2} - \sqrt{(d+L)^2 + R^2} \right) $$
\end{multicols*}
\end{document}